AI agents write their own long-term memory, retrieve it by semantic search, and keep it on the
user's device, where encryption at rest is the default defense. This paper studies an adversary
that obtains the storage layer partway through the life of such a store, for example through a
compromised synchronization account or repeated disk imaging, observes later activity, and can
steer the agent's inputs, but holds only stored state for records written earlier. For encrypted
vector stores whose layout follows content, labels learned from induced writes transfer to the
earlier records, and a transfer lemma shows the two cohorts are then exposed equally. The lemma
also names its own precondition: transfer requires that the layout disclose which partition holds a
given earlier record, which a growing layout does and a sealed fixed-size one does not. Subject
drift appears to protect early records, but the adversary chooses what it induces, and one that
induces with archival content rather than current content recovers the advantage in full, from
$\REfourNaive$ to $\REfourAware$ on a timestamped corpus. Memories a language-model agent actually wrote carry
the structure the attack needs.

Against a store that withholds that precondition, three channels supply it anyway. A sequence of
retained file versions reconstructs pre-compromise writes, contradicting the claim that history
independence makes several snapshots no worse than one; the region holding an earlier record lies
inside the version diff at every capture granularity tested. A bounded per-partition capacity
discloses the earlier cohort's load whenever a write is refused, at near-perfect accuracy, unless a
refusal performs the writes an acceptance would have. And an agent that acts on what it recalls
leaks through the mere occurrence of a write, which makes deduplication a membership oracle for
facts written before the compromise. Four of nine open-source agent-memory frameworks, read from
their shipped source, make whether a write happens depend on what is already stored.

We define dual-epoch security as an adaptive, leakage-parameterized game, give sufficient conditions
on the storage layout, and show that forward and backward privacy do not imply it. We then evaluate
three designs. A memory-resident store with content-oblivious placement and clocked flushes closes
all three channels, because its writes carry no content, its searches touch no disk, and a flush
occurs whether or not the agent wrote. Exhaustive search in memory is faster than the encrypted
on-disk alternative only up to about $10^4$ records, but the partition index that costs a
content-placed store its write channel costs this one nothing when it is kept in locked memory
rather than on disk, and restores its lead at every size measured. A partitioned sealed store
covers larger corpora with its residual channels stated, including one that is easy to
misattribute: its keyed partition permutation is largely recovered from ordinary search traffic by
an adversary that observes reads.
